% ============================================================================
% Core Feature: Access Control (ACL)
% ============================================================================
\section{Access Control (ACL)}

\begin{ugModule}{Overview}

  \textbf{Purpose:}
  Manage fine-grained, per-document access permissions with
  multi-type grants and a self-service access request workflow.

  \textbf{Who Can Access:}
  \ugRole{Admin} \ugRole{Manager} (manage);
  \ugRole{Staff} (request access)
\end{ugModule}

% ────────────────────────────────────────────────────────────────
\subsection{Access Control Model}

InsightDoc uses a three-tier access control model. The
document-level ACL is set during document creation (Tab~3)
and can be changed later:

\begin{tabularx}{\textwidth}{@{} l X @{}}
  \toprule
  \textbf{ACL Level} & \textbf{Behavior} \\
  \midrule
  \textbf{Public}
    & All authenticated users can access. No additional
      configuration required. \\
  \textbf{Protected}
    & All authenticated users within the organization can
      access. External sharing is restricted. \\
  \textbf{Private}
    & Only explicitly granted users, roles, or units can
      access. The document is invisible to everyone else. \\
  \bottomrule
\end{tabularx}

For Private documents, access is granted through
\texttt{DocumentAccess} records that specify exactly who
can view and download the document.

% ────────────────────────────────────────────────────────────────
\subsection{Access Grant Types}

When a document is set to Private, access can be granted
through three different types:

\begin{tabularx}{\textwidth}{@{} l l X @{}}
  \toprule
  \textbf{Type} & \textbf{Code}
    & \textbf{Description} \\
  \midrule
  By Unit
    & \texttt{unit}
    & Grants access to \textbf{all users} within a
      selected organizational unit. When new users
      are added to the unit, they automatically
      receive access. \\
  By Role (Division)
    & \texttt{role}
    & Grants access by unit division level
      (positional role). All users holding that
      position receive access. \\
  By User
    & \texttt{user}
    & Grants access to a \textbf{specific individual
      user}. Most granular control. \\
  \bottomrule
\end{tabularx}

% ────────────────────────────────────────────────────────────────
\subsection{Access Grant Data}

Each \texttt{DocumentAccess} record stores:

\begin{tabularx}{\textwidth}{@{} l X @{}}
  \toprule
  \textbf{Field} & \textbf{Description} \\
  \midrule
  Document
    & The document this grant applies to \\
  Access Type
    & Unit, Role (Division), or User \\
  Unit
    & Target unit (when type is ``unit'') \\
  Unit Division
    & Target division role (when type is ``role'') \\
  Allow User
    & Target individual user (when type is ``user'') \\
  Granted By
    & Name of the person who created this grant \\
  Granted At
    & Timestamp when access was granted \\
  Expires At
    & Optional expiry date; access is automatically
      denied after this date \\
  \bottomrule
\end{tabularx}

\begin{ugTip}
  Use time-limited access grants (\ugField{Expires At}) for
  temporary collaborators, contractors, or audit purposes.
  This eliminates the need to manually revoke access when
  the collaboration period ends.
\end{ugTip}

% ────────────────────────────────────────────────────────────────
\subsection{Managing Access}

Access is managed via the document's \ugMenu{Access} tab:

\begin{tabularx}{\textwidth}{@{} l l X @{}}
  \toprule
  \textbf{Action} & \textbf{API Endpoint}
    & \textbf{Description} \\
  \midrule
  View grants
    & \texttt{GET .../access}
    & List all current access grants \\
  Add grant
    & \texttt{POST .../access}
    & Grant access to a unit, role, or user \\
  Revoke grant
    & \texttt{DELETE .../access/\{id\}}
    & Immediately revoke a specific grant \\
  \bottomrule
\end{tabularx}

% ────────────────────────────────────────────────────────────────
\subsection{Access Request Workflow}

Users who lack access to a Private document can submit a
self-service access request:

\subsubsection*{Request Fields}

\begin{tabularx}{\textwidth}{@{} l X @{}}
  \toprule
  \textbf{Field} & \textbf{Description} \\
  \midrule
  Request Code
    & Auto-generated unique request identifier \\
  Document
    & The document being requested \\
  Access Type
    & The type of access requested \\
  Requested User
    & The user submitting the request \\
  Requested By
    & Name of the requesting user \\
  Requested At
    & Timestamp of submission \\
  Requested Message
    & Mandatory reason/justification for the
      request \\
  Status
    & Current status: Pending, Approved, or
      Rejected \\
  \bottomrule
\end{tabularx}

\subsubsection*{Response Fields (After Review)}

\begin{tabularx}{\textwidth}{@{} l X @{}}
  \toprule
  \textbf{Field} & \textbf{Description} \\
  \midrule
  Responded User
    & The admin/manager who reviewed the request \\
  Responded By
    & Name of the reviewer \\
  Responded At
    & Timestamp of the review decision \\
  Response Comment
    & Optional comment from the reviewer \\
  Document Access
    & If approved, the resulting
      \texttt{DocumentAccess} record \\
  Expires At
    & Optional expiry set by the reviewer \\
  \bottomrule
\end{tabularx}

% ────────────────────────────────────────────────────────────────
\subsection{Step-by-Step: Requesting Access}

\begin{ugSteps}
  \item Navigate to a document you cannot access. The system
        shows a restricted access message.
  \item Click \ugButton{Request Access}.
  \item Enter a \textbf{justification message} explaining
        why you need access.
  \item Click \ugButton{Submit Request}.
  \item Your request enters \texttt{Pending} status.
  \item You will be notified when the request is approved
        or rejected.
\end{ugSteps}

\subsection{Step-by-Step: Reviewing Access Requests}

\begin{ugSteps}
  \item Navigate to the document's detail page.
  \item Open the \ugMenu{Access Requests} section.
  \item Review pending requests --- see the requester's name,
        message, and submission date.
  \item Click \ugButton{Approve} to grant access or
        \ugButton{Reject} to deny.
  \item If approving, a \texttt{DocumentAccess} entry is
        automatically created, linking the request to the
        new grant.
  \item Optionally set an \ugField{Expires At} date for the
        grant.
  \item Add a \ugField{Response Comment} if needed.
\end{ugSteps}

% ────────────────────────────────────────────────────────────────
\subsection{Activity Tracking}

All access-related actions are recorded as
\texttt{DocumentActivity} entries:

\begin{tabularx}{\textwidth}{@{} l X @{}}
  \toprule
  \textbf{Action Code} & \textbf{Description} \\
  \midrule
  \texttt{request-access}
    & User submitted an access request \\
  \texttt{approve-access}
    & Admin/Manager approved a request \\
  \texttt{reject-access}
    & Admin/Manager rejected a request \\
  \bottomrule
\end{tabularx}

% ────────────────────────────────────────────────────────────────
\subsection{Business Rules}

\begin{itemize}[leftmargin=18pt]
  \item Access grants are enforced during document preview,
        download, and detail view operations.
  \item Expired access grants (\texttt{ExpiresAt} in the past)
        are automatically denied.
  \item Only one pending request per user per document is
        allowed (duplicate prevention).
  \item Approved requests automatically create a linked
        \texttt{DocumentAccess} record.
  \item Revoking an access grant is immediate and permanent.
  \item Unit-based grants are dynamic --- when a new user
        joins the unit, they automatically receive access.
  \item Role/Division grants work similarly --- all users
        holding the specified position receive access.
\end{itemize}

\begin{ugImportant}
  When a document is set to \textbf{Private} during creation,
  at least one access grant (by unit, role, or user) must be
  specified. This prevents creating orphan documents that no
  one can access.
\end{ugImportant}

\newpage
